\section{Results and discussion}

The results are presented in the order of the audit. The first two experiments
test whether the method rejects known or plausible non-identifiable systems.
Sensorium then establishes that predictive utility can coexist with a negative
mechanistic decision. MICRONS finally tests whether a positive local
structure-function result survives matched controls and independent cohorts.

\subsection{MIS recovers synthetic truth and rejects the Allen mechanism}

The known-truth experiment behaved as specified. The directed system passed all
three evidence blocks and obtained (M=1.0). The common-drive system passed
reproducibility but failed topology specificity and directionality, yielding
(M=0.333). The topology-only system passed reproducibility and topology but
failed directionality ((M=0.667)). The direction-only system failed topology
specificity despite temporal information ((M=0.719)). Thus, only the case with
the complete designed mechanism passed the conjunctive decision.

These cases illustrate why the binary decision cannot be replaced by the
descriptive score. The direction-only system has a larger (M) than the
topology-only system, but neither supports the complete claim. Their scores
describe different patterns of partial evidence; they do not provide a ranking
of mechanistic truth. The experiment also confirms that high reproducibility
does not rescue a common-drive explanation.

Allen produced the intended real negative control. Median cross-mouse
correlation was 0.891, median split-half correlation was 0.990, and 75\% of
instances exceeded their own 95th-percentile null. The reproducibility block
therefore passed. In contrast, the Allen-informed topology performed 0.0054
worse than the median permutation, outperformed only 21\% of permutations, and
had a paired bootstrap lower bound of (-0.0075). Its advantage over a
disconnected model was only 0.0035, below the fixed 0.05 threshold.

Directed identifiability also failed: the evaluated latency and lead-lag
fractions were zero under the operational criteria. The final MIS was 0.358 and
the decision was \texttt{reproducible\_target\_without\_mechanistic\_
identifiability}. This does not imply that Allen responses lack biological
structure. It means that the current target and model comparison do not identify
the proposed directed anatomical mechanism. That distinction is the primary
purpose of the negative control.

\subsection{Sensorium supports prediction but not mechanistic identification}

Static Sensorium showed consistent predictive utility. Across seven validation
mice, the median best correlation was 0.328, the median improvement over mean
response was 0.102, and the improvement over scrambled stimulus was 0.077.
Reliability cannot be estimated from non-repeated validation stimuli and is
therefore not assigned a favorable value. In the five repeated-test mice, median
reliability was 0.642 and median best correlation was 0.341. Improvements over
mean response and scrambling were 0.095 and 0.068, respectively.

The separate topographic constraint passed in all five repeated-test datasets,
with median Spearman association 0.177 and median effect over its null 0.178.
This is partial structural evidence, but it does not identify synaptic direction
or a causal circuit. Consequently, no static Sensorium mouse passed the complete
MIS. The scientifically appropriate conclusion is that stimulus and context
contain repeatable predictive information and a modest topographic signal, not
that the fitted predictors recover the generating mechanism.

Dynamic Sensorium produced best per-mouse correlations from 0.427 to 0.526 in
the current public cohort. The temporal filterbank won one mouse, temporal SVD
two, and the compact multilayer perceptron two. Relative to mean response, the
temporal filterbank improved all five mice with a median gain of 0.031. These
comparisons show that temporal modelling adds predictive information, but the
small number of mice and modest model differences do not justify a broad model
superiority claim.

The bounded official-architecture run completed through the official loader and
model factory for five mice, but correlations remained between 0.011 and 0.035
under the deliberately restricted training budget. It is therefore reported as
an integration check only. Presenting this artifact as an official Sensorium
baseline or a state-of-the-art comparison would be misleading. The audit keeps
it visible while blocking that interpretation.

\subsection{MICRONS yields a replicated local structure-function association}

The three MICRONS cohorts jointly contain 2991 units, 6575 loaded synapses, and
5943 unique connected directed pairs. The discovery, first hold-out, and second
hold-out contain 2095, 1926, and 1922 connected pairs among 999000, 983072, and
997002 candidate directed non-self pairs. The primary endpoint passed the fixed
distance- and degree-matched tests in every cohort.

\begin{table}[t]
\centering
\caption{Replicated MICRONS primary endpoint: connected directed pairs show closer readout-location similarity than matched non-connected controls.}
\label{tab:microns-primary}
\begin{tabular}{lrrrrrrrr}
\toprule
Cohort & Units & Synapses & Edges & Tests & $\Delta_d$ & $q_d$ & $\Delta_k$ & $q_k$ \\
\midrule
Discovery & 1000 & 2267 & 2095 & 28 & 0.0215 & 0.0076 & 0.0403 & 0.0053 \\
Hold-out 1000 & 992 & 2161 & 1926 & 30 & 0.0210 & 0.0070 & 0.0380 & 0.0062 \\
Hold-out 2000 & 999 & 2147 & 1922 & 25 & 0.0179 & 0.0087 & 0.0286 & 0.0083 \\
\bottomrule
\end{tabular}
\end{table}


For distance-matched controls, the connected-pair advantages in readout-location
similarity were 0.0215, 0.0210, and 0.0179, with one-sided false-discovery-rate
adjusted (q)-values of 0.0076, 0.0070, and 0.0087. Degree-matched effects were
larger: 0.0403, 0.0380, and 0.0286, with adjusted (q)-values of 0.0053, 0.0062,
and 0.0083. The effect therefore retained its sign and approximate magnitude
outside the discovery window.

\begin{table}[t]
\centering
\caption{Unit-cluster bootstrap stability analysis for the MICRONS primary endpoint using 300 resamples per cohort. PI denotes a 95\% percentile interval. Intervals are reported as robustness summaries for neuron-sharing pair data, not as naive pair-level uncertainty estimates.}
\label{tab:microns-bootstrap}
\small
\begin{tabular}{lrrrr}
\toprule
Cohort & Distance median & Distance 95\% PI & Degree median & Degree 95\% PI \\
\midrule
Discovery & 0.0218 & [0.0166, 0.0263] & 0.0405 & [0.0349, 0.0454] \\
Hold-out 1000 & 0.0209 & [0.0142, 0.0263] & 0.0377 & [0.0321, 0.0431] \\
Hold-out 2000 & 0.0181 & [0.0115, 0.0238] & 0.0287 & [0.0229, 0.0354] \\
\bottomrule
\end{tabular}
\end{table}


Unit-cluster bootstrap intervals remained strictly positive. The
distance-matched lower bounds were 0.0166, 0.0142, and 0.0115; the
degree-matched lower bounds were 0.0349, 0.0321, and 0.0229. Resampling units
rather than pairs reduces the false precision that would result from treating
all pairwise observations as independent. Together with matched controls and
hold-out replication, these intervals support a robust local association.

The result is narrower than a causal mechanism. Readout-location similarity is
a functional descriptor associated with connectivity after two important
controls, but unmeasured cell type, laminar organization, or other local
variables may contribute. The analysis does not perturb edges, establish
necessity, or predict behaviour. Its value is that the admitted claim is both
positive and specific: within the evaluated MICRONS volume, connected directed
pairs are more similar in readout location than comparable non-connected pairs.

This result is also not a new biological wiring rule. Ding et al. have already
reported a substantially more detailed MICRONS analysis of like-to-like
connectivity using receptive-field, feature, and axon--dendrite geometry
controls \citep{ding2025wiring}. Agreement with that work provides an
external-validity check for the MouseBrainBench pipeline. The potential novelty
here is the same claim-control protocol operating across synthetic, Allen,
Sensorium, and MICRONS regimes, not priority over the underlying
structure-function phenomenon.

\subsection{Cross-regime interpretation and limitations}

The experiments demonstrate three distinct outcomes that a correlation-only
evaluation would blur. Allen supplies reliable neural targets without
identifiable topology. Sensorium supplies useful in-distribution and temporal
prediction without anatomical or causal identification. MICRONS supplies a
replicated local structure-function association after matched anatomical
controls. MouseBrainBench does not force these outcomes onto a single scale; it
assigns each one the strongest interpretation that its evidence permits.

Sensitivity analyses preserve the main qualitative decisions under reasonable
threshold changes: Allen remains a negative mechanistic case, Sensorium remains
predictive evidence, and MICRONS retains the preregistered local endpoint across
two hold-outs. This supports the methodological claim that the conclusions are
not consequences of one arbitrary threshold. It does not establish that the
current thresholds are universal; new targets require domain-specific criteria
declared before evaluation.

Several limitations remain material. The positive MICRONS result is local,
observational, and dominated by one functional descriptor. The Allen analysis
uses a controlled subset rather than every available session. Dynamic Sensorium
contains only five current and five legacy mice in the present local benchmark,
and the official architecture was not trained with a leaderboard-equivalent
budget. Finally, MIS currently encodes necessary conditions selected for these
claims; it is not a proven complete theory of mechanistic explanation.

These limits determine the publication claim. The framework and experiments
support a reproducible audit of partial mouse-brain models and a replicated
local MICRONS association. They do not support a complete mouse-brain digital
twin, consciousness, behavioural equivalence, causal synaptic mechanisms, or a
state-of-the-art Sensorium predictor. Keeping those statements blocked is part
of the reported result rather than an omission.
